\subsection{Recepción continua seleccionada de E2}
Synaptix selecciona la recepción continua para cumplir producción E2 y Operación de ambos alcances desde el inicio del mes 21. Para inicio el 1 de diciembre de 2026, la fecha es el 1 de agosto de 2028. La línea base adopta como supuesto de planificación la participación de la Contraparte en la revisión continua y en la sesión final de aceptación, con firma expresa prevista al cierre de la evidencia. Synaptix propone este procedimiento para su ejecución contractual; no acredita un acuerdo ni un acta ya suscritos.~\parencite[arts.~17.2.4, 17.3 y 18; pp.~12--13]{bases-admin}

El expediente parcial se prepara desde el 15 de junio, con índice, versión, conciliación, formación, fuentes y fórmulas SLA, incidentes y observaciones. La red reserva diez días hábiles de revisión parcial y diez de respuesta sobre evidencia ya producida. Calidad entregará registros diarios inmutables y resolverá observaciones tempranas; la revisión parcial no acepta mediciones futuras ni consume el derecho de revisar evidencia nueva.

Las cuatro últimas semanas se observan íntegramente del 4 al 31 de julio, con volumen real, SLA y usuarios. El 31 de julio, de 20:00 a 24:00, se reservan dos ejecutores propios separados, uno de Calidad y otro de Proyecto, cuatro HH cada uno. El cuadrante conserva el intervalo natural, sin reutilizar un turno NOC/SOC/mesa o una jornada de ingeniería. Revisan el expediente vivo y la última evidencia; el cierre automático a las 00:00 conserva el registro final e impide habilitación ante diferencias, incidentes impeditivos o medición incompleta. Ninguna proyección sustituye los últimos datos.

El acta expresa de cierre de marcha blanca es una puerta previa del corte. Synaptix preparará y ensayará en PREPROD la configuración, respaldo, retorno y cambio de autoridad antes de la ventana; el corte es una habilitación lógica preparada, cuya secuencia será cierre de evidencia, comprobación final, acta expresa y activación. El ensayo medirá los tiempos de consolidación y activación: una precedencia lógica no acredita duración física nula ni firma instantánea. Sólo se habilitará E2 como registro oficial si se dispone del acta y de las condiciones copulativas. NOC, SOC y mesa atienden ambos alcances desde el inicio de 21; la estabilización de 28 días comienza el 1 de agosto, con 448 HH y cuadrante propio.

H12 corresponde a la aceptación final del proyecto de implementación dentro del mes 21 y agrega las entregas finales, incluido el cierre de pilotos de innovación y sus revisiones. Se distingue del acta de marcha blanca requerida para el corte del primer día. Esa distinción no autoriza producción sin acta ni aceptación anticipada de un producto.

\paragraph{Preparación y autorización del cierre.} Antes de H10, Proyecto presentará a la Contraparte la agenda, referente y suplente, criterios, plantilla de acta, expediente diario y protocolo de cierre para coordinación expresa. La convocatoria y la respuesta se registrarán, sin presumir disponibilidad ni aceptación por silencio. Antes de H11, Calidad e Integración ensayarán en PREPROD la consolidación del último registro, verificación, presentación del acta, habilitación y retorno; conservarán duración medida de cada paso, versión y responsable. El ensayo no incluirá una firma simulada como evidencia de decisión del CLIENTE ni convertirá tiempo de elaboración en tiempo de aceptación externa.

Durante los 28 días finales, el tablero diario contrastará mínimos de volumen, cobertura, SLA, conciliación, formación e incidentes; Proyecto escalará una previsión de incumplimiento tan pronto sea identificada, sin esperar al último día. A siete días del cierre comprobará agenda y suplentes, observaciones abiertas y preparación del corte. El acta sólo se someterá a firma cuando esté completa la última evidencia. Una duración positiva medida, una observación o un pronunciamiento tardío se reflejará en la programación de la transición: el margen cero no absorbe esa demora. Sin acta expresa y comprobación final no se habilita el registro oficial. Este protocolo se incorpora al expediente de aceptación de T-17; H12 posterior no reemplaza la autorización previa.\parencite[arts.~17--18, pp.~12--13]{bases-admin}

\paragraph{Supuesto adoptado y contingencia.} Para la línea base ofrecida se supone que la Contraparte participará en la sesión final del 31 de julio de 2028, de 20:00 a 24:00, y emitirá su pronunciamiento expreso una vez cerrado el registro de las cuatro semanas. Proyecto coordinará la sesión y presentará el acta; Calidad comprobará la evidencia completa. La firma y activación se programan en la transición al 1 de agosto, conservando ese orden y todos los criterios. Proyecto presentará con la línea base la agenda de recepción continua y el protocolo de cierre; el CLIENTE resolverá su coordinación mediante un pronunciamiento expreso. La viabilidad del instante de corte depende de ese procedimiento y de los tiempos medidos en el ensayo; la reserva cero exige tratar cualquier demora como desviación, no como holgura disponible. La decisión de planificación queda cerrada bajo este supuesto, sin atribuir una aceptación ya obtenida. El CLIENTE conserva los diez días de revisión de evidencia nueva: el supuesto prevé un pronunciamiento anticipado válido, sin renuncia a ese derecho ni aceptación tácita. Si falta el acta, existen observaciones impeditivas o se usa íntegramente la revisión posterior, la puerta bloquea el cambio oficial, se registra el atraso y se extiende marcha blanca a costo de Synaptix, sin desplazar el calendario contractual. PL-R04 controla el incumplimiento del supuesto; el escenario máximo se conserva como sensibilidad de plazo.
